\subsection{Constructing the geometric prior family}

Let $Y-X-Z$ and let the encoder posterior be
$q_\phi(z|x)=\mathcal N(\mu_\phi(x),\operatorname{diag}\sigma_\phi^2(x))$.
CEB~\citep{fischer2020conditional} replaces the VIB marginal reference by a
label-conditional reference $r(z|y)$ and uses the certificate
\begin{equation}
\E\KL(q_\phi(z|X)\|r(z|Y))
=I(X;Z|Y)+\E_Y\KL(q_\phi(z|Y)\|r(z|Y))\ge I(X;Z|Y).
\label{eq:gap}
\end{equation}
At the class-wise matched optimum $r=q(z|y)$, the certificate reduces to
\begin{equation}
\mathrm{ReX}=I(X;Z|Y),
\label{eq:reix-residual}
\end{equation}
and the certificate contains no term coupling two distinct label conditionals.
Consequently, it records neither their relative directions nor their pairwise
distances, so the compression criterion alone cannot prefer any cross-label geometry.

\paragraph{The GPB framework.}
GPB is not a single prior but a framework for declaring reference geometry: it
restricts the reference to stateless conditional Gaussians
\begin{equation}
p_\theta(z|y)=\mathcal N(m_\theta(y),\Sigma_\theta(y)),\qquad
\Sigma_\theta(y)\preceq\tau_{\max}^2I,
\label{eq:pgib-covariance}
\end{equation}
whose means obey a declared two-sided separation for every pair of distinct
discrete labels $y\ne y'$,
\begin{equation}
\Delta\le\|m_\theta(y)-m_\theta(y')\|\le\bar\Delta,
\label{eq:pgib-separation-condition}
\end{equation}
or a two-sided label metric condition for $P_Y$-almost every label pair
$(y,y')$, where $d_Y$ is a metric on the label space $\mathcal Y$:
\begin{equation}
\rho_{\min}d_Y(y,y')\le\|m_\theta(y)-m_\theta(y')\|
\le\rho_{\max}d_Y(y,y'),
\label{eq:pgib-lipschitz}
\end{equation}
with constants fixed before training. These conditions define the geometric prior
family $\mathcal G$. Any method whose reference prior belongs to $\mathcal G$ and
whose readout is tied to the same anchor means enters the framework and inherits the
guarantees of Section~\ref{sec:guarantees}; OPB (classification) and EPB (regression),
studied below, are two members, not the extent of the framework. The objective is
\begin{equation}
\mathcal L_{\rm GPB}=\E[\ell(c_\psi(z),Y)+\beta\KL(q_\phi(z|X)\|p_\theta(z|Y))],
\qquad z\sim q_\phi(z|X).
\label{eq:pgib-objective}
\end{equation}
The instances below realize this single objective: the anchor means enter the KL and
are reused as prediction anchors through the tied readout (anchor-energy scores for
classification, tied projection for regression), so the declared geometry is shared by
compression and prediction and no instance-specific objective is stated. Here $\ell\ge0$
is cross-entropy for classification and the chosen nonnegative regression loss for
continuous labels.

Because each GPB prior $p_\theta(z|y)$ depends only on $y$, the CEB certificate
identity applies directly with $r=p_\theta$. Restricting the conditional reference
to the geometric prior family therefore preserves the certificate and its
lower-bound interpretation, while changing the cross-label geometry represented by
the feasible prior family.

\paragraph{Instances.}\label{sec:opb}
For classification, OPB evaluates a prior encoder on all class labels and uses its
thin polar factor $Q_\theta$. If $U_\theta$ is the resulting all-class matrix, then
\begin{equation}
Q_\theta=U_\theta(U_\theta^\top U_\theta)^{-1/2},
\qquad Q_\theta^\top Q_\theta=I_K,
\label{eq:opb-qr}
\end{equation}
and we assume, monitoring rather than imposing it, that the monitored matrix has
full column rank (Assumption~\ref{ass:fullrank} in the appendix). Writing
$q_{\theta,k}$ for the $k$-th column of the frame $Q_\theta$, the anchors are
\begin{equation}
p_\theta^{\rm OPB}(z|Y=k)=\mathcal N(aq_{\theta,k},\tau^2I),\qquad
Q_\theta^\top Q_\theta=I_K,\qquad\|aq_i-aq_j\|=\sqrt2a.
\label{eq:opb-prior}
\end{equation}
The construction also gives the fixed geometry
\begin{equation}
(aq_i)^\top(aq_j)=a^2\mathbf 1\{i=j\},
\qquad \|aq_i-aq_j\|=\sqrt2a\quad(i\ne j).
\label{eq:opb-geometry}
\end{equation}
The anchor-energy scores are
\begin{equation}
s_k(z)=-\frac{\|z-aq_{\theta,k}\|^2}{2\tau^2}+\log\pi_k,
\label{eq:opb-energy}
\end{equation}
the Gaussian Bayes scores; with uniform priors this induces an orthogonal classifier
whose scale is fixed by the prior (Appendix~\ref{app:construction}). Detailed polar
construction, initialization, covariance relaxation, and the relation to
prototype/ETF methods are given in Appendices~\ref{app:construction} and~\ref{sec:related}.

For regression, EPB uses a one-dimensional continuous label, standardized once on
the training set as $\tilde y$ with the statistics frozen. Let $W\in\R^{d\times1}$ be
a trainable direction with per-step normalization $u=W/\|W\|$; the isometric prior is
\begin{equation}
p_\theta^{\rm EPB}(z|y)=\mathcal N(\rho\tilde y\,u,\tau^2I),\qquad
\|\mu_p(y)-\mu_p(y')\|=\rho|\tilde y-\tilde y'|,
\label{eq:epb-prior}
\end{equation}
so latent distances equal $\rho$ times standardized label distances exactly. Order
and magnitude of label differences are encoded with equality, which makes it a
member of $\mathcal G$ (Eq.~\eqref{eq:pgib-lipschitz} with
$\rho_{\min}=\rho_{\max}=\rho$). The
prediction head is the tied projection $\hat{\tilde y}=u^\top z/\rho$, the regression
analogue of the anchor-energy classifier: no free direction or scale can bypass the
prior. The main variant uses the trainable normalized axis with no bias; the
fixed-axis and unconstrained-scale ablations are in
Appendix~\ref{app:construction}.

\subsection{Comparator bounds and posterior-center separation}
\label{sec:guarantees}
Define
\begin{equation}
\mathcal R(w)=\E[\ell(c_\psi(Z),Y)],\qquad
\mathcal K(w)=\E\KL(q_\phi(z|X)\|p_\theta(z|Y)),
\label{eq:risk-rate-definitions}
\end{equation}
and
\begin{equation}
\mathcal D(w)=\frac{\E\|\mu_\phi(X)-m_\theta(Y)\|^2}{2\tau_{\max}^2}.
\label{eq:expected-anchor-mismatch}
\end{equation}
\begin{proposition}[Comparator-dependent bound on anchor mismatch]
\label{prop:pgib-alignment}
Let $\beta>0$, $\varepsilon\ge0$, and $\ell\ge0$. If $\widehat w$ is
$\varepsilon$-optimal for $\mathcal R+\beta\mathcal K$, then for every feasible
comparator $w_0$ in the same family (same encoder, posterior, prior, and
prediction-head parameterization),
\begin{equation}
\mathcal D(\widehat w)\le\mathcal K(\widehat w)\le
\mathcal K(w_0)+\frac{\mathcal R(w_0)+\varepsilon}{\beta}.
\label{eq:pgib-comparator-bound}
\end{equation}
Specializing to a comparator with total KL $\delta_{\rm app}$ and risk
$C_{\rm app}$ gives
\begin{equation}
\mathcal D(\widehat w)\le\delta_{\rm app}
 +\frac{C_{\rm app}+\varepsilon}{\beta}.
\label{eq:pgib-approximate-alignment}
\end{equation}
With a near-realizable comparator, namely small $\delta_{\rm app}$,
Eq.~\eqref{eq:pgib-approximate-alignment} becomes an approximate-realizability
statement.
\end{proposition}

The bound certifies alignment to whatever comparator is prespecified: its KL is the
comparator-specific certificate floor, and the risk term controls how fast the
right-hand certificate bound approaches that floor as $\beta$ grows. It does not
claim that the realized mismatch $\mathcal D(\widehat w_\beta)$ decreases
monotonically in $\beta$, converges to, or attains that floor: with a fixed
comparator and a fixed error bound, only the right-hand certificate approaches
$\mathcal K(w_0)$, which is a property of the chosen comparator rather than an
intrinsic dataset quantity.

For a target $d_\star$, the bound is informative only if
$\mathcal K(w_0)<d_\star$ and
\begin{equation}
\beta\ge\frac{\mathcal R(w_0)+\varepsilon}{d_\star-\mathcal K(w_0)}.
\label{eq:comparator-informative-condition}
\end{equation}
A prespecified comparator's KL and risk can be estimated on an independent held-out
set using
\begin{equation}
\widehat{\mathcal K}_S(w_0)=\frac1{|S|}\sum_{(x_i,y_i)\in S}
\KL(q_{\phi_0}(z|x_i)\|p_{\theta_0}(z|y_i)),\quad
\widehat{\mathcal R}_S(w_0)=\frac1{|S|}\sum_{(x_i,y_i)\in S}
\E_{q_{\phi_0}}\ell(c_{\psi_0}(Z),y_i),
\label{eq:heldout-comparator-estimates}
\end{equation}
which estimates a comparator-specific floor, not an intrinsic dataset floor.

\begin{theorem}[Mismatch-conditioned posterior-center separation]
\label{thm:pgib-separation}
For discrete labels, suppose $\E\|\mu_\phi(X)\|^2<\infty$, let
$m_k=\E[\mu_\phi(X)\mid Y=k]$ and
\begin{equation}
D_k=\frac{\E[\|\mu_\phi(X)-m_\theta(k)\|^2\mid Y=k]}{2\tau_{\max}^2},
\qquad
r_k=\sqrt{2\tau_{\max}^2D_k}.
\label{eq:pgib-class-rate}
\end{equation}
Then
\begin{equation}
\|m_i-m_j\|
\ge\|m_\theta(i)-m_\theta(j)\|-r_i-r_j
\ge\Delta-r_i-r_j.
\label{eq:pgib-separation}
\end{equation}
For continuous labels, suppose $\E\|\mu_\phi(X)\|^2<\infty$, define
$m(y)=\E[\mu_\phi(X)\mid Y=y]$ and
\begin{equation}
D(y)=\frac{\E[\|\mu_\phi(X)-m_\theta(y)\|^2\mid Y=y]}{2\tau_{\max}^2},
\qquad
r(y)=\sqrt{2\tau_{\max}^2D(y)},
\label{eq:pgib-class-rate-continuous}
\end{equation}
and the analogous two-level lower bound is
\begin{equation}
\|m(y)-m(y')\|
\ge\|m_\theta(y)-m_\theta(y')\|-r(y)-r(y')
\ge\rho_{\min}d_Y(y,y')-r(y)-r(y'),
\label{eq:pgib-separation-continuous}
\end{equation}
for $P_Y\otimes P_Y$-almost every pair $(y,y')$. The first-level bound in
Eq.~\eqref{eq:pgib-separation} is positive exactly when
\begin{equation}
\Gamma_{ij}:=\frac{r_i+r_j}
{\|m_\theta(i)-m_\theta(j)\|}<1,
\label{eq:separation-tightness-ratio}
\end{equation}
while the declared-margin level is positive exactly when $r_i+r_j<\Delta$. The two
conditions coincide for OPB, whose anchor pair distance is the constant $\sqrt2a$ of
Eq.~\eqref{eq:opb-geometry}, so the declared margin satisfies $\Delta=\sqrt2a$; with
equal mismatch $D_i=D_j=D$ the common condition reduces to
$D<a^2/(4\tau_{\max}^2)$. In the continuous case, the declared metric lower bound is
positive exactly when
\begin{equation}
r(y)+r(y')<\rho_{\min}d_Y(y,y'),
\label{eq:separation-tightness-ratio-continuous}
\end{equation}
so nearby label pairs admit vacuous bounds more readily. The scale is a
validation-selected design parameter and does not itself guarantee a positive bound.
The declared margin thus survives up to the two measured correction terms, so the
guarantee is informative exactly when the measured mismatch stays below the margin;
small measured mismatch transfers the declared geometry to the posterior centers
almost intact.
\end{theorem}

Combining the objective-to-mismatch proposition with the posterior-center
separation theorem gives the main end-to-end guarantee of this section.
\begin{corollary}[Beta-dependent transfer of anchor separation]
\label{cor:pgib-beta-transfer}
Fix the prior family, a feasible comparator $w_0$, and an
$\varepsilon_\beta$-optimal solution $\widehat w_\beta$ for each $\beta>0$, with
$\sup_\beta\varepsilon_\beta<\infty$. Define
\begin{equation}
B_\beta:=\mathcal K(w_0)+\frac{\mathcal R(w_0)+\varepsilon_\beta}{\beta}.
\label{eq:pgib-beta-budget}
\end{equation}
For discrete labels, suppose $\pi_k\ge\pi_{\min}>0$ for all considered classes. Then
for every pair $i\ne j$,
\begin{equation}
\|m_i-m_j\|
\ge \Delta-2\sqrt{\frac{2\tau_{\max}^2 B_\beta}{\pi_{\min}}}.
\label{eq:pgib-beta-separation}
\end{equation}
For continuous labels, for any $\xi\in(0,1)$ and $B_\beta>0$, there is a set
$S_{\beta,\xi}$ with $P_Y(S_{\beta,\xi})\ge 1-\xi$ such that, for
$P_Y\otimes P_Y$-almost every $(y,y')\in S_{\beta,\xi}^2$,
\begin{equation}
\|m(y)-m(y')\|
\ge \rho_{\min}d_Y(y,y')
-2\sqrt{\frac{2\tau_{\max}^2 B_\beta}{\xi}}.
\label{eq:pgib-beta-separation-continuous}
\end{equation}
With a fixed comparator satisfying $\mathcal K(w_0)=0$ and the remaining quantities
bounded while $\beta$ varies, the certified separation deficit is
$O(\beta^{-1/2})$; with $\mathcal K(w_0)>0$, the certificate approaches the
comparator-dependent non-zero upper bound $\mathcal K(w_0)$. These rate statements
concern the certificate only: they do not claim that the realized mismatch
$\mathcal D(\widehat w_\beta)$ or the realized center distances converge at this
rate.
\end{corollary}
